% GENERATED FILE. Edit canonical lessons, then run npm run generate:books.
% canonical-source-hash: fnv1a64:0fb5752ed0de4ad9
% canonical-lessons: UR-C399-hifazat, UR-C399-talashi, UR-C399-kastam, UR-C399-mamnu, UR-C399-charhana

\chapter{Security Checks}
\label{ch:ur-security-checks}
\hlchaptermodality{399}

\begin{chapteropening}
\textbf{By the end of this chapter:} \emph{I can get through a security check and customs.}

Complete the last lesson of chapter 399: I can get through a security check and customs.
\end{chapteropening}

\section[\texorpdfstring{\ur{حفاظت} (hifāzat)}{حفاظت (hifāzat)}]{\ur{حفاظت} (hifāzat) — protection; safety}

\label{lesson:UR-C399-hifazat}



\begin{quote}
Before the new one: say the Urdu for to tighten, then the Urdu for postponed.
\end{quote}

\subsection*{You'll want to know: \ur{حفاظت}}
\textbf{\ur{حفاظت}} — \emph{hifāzat} — \textquotedblleft{}protection; safety\textquotedblright{}.

\begin{quote}
\ur{مسافروں} \ur{کی} \ur{حفاظت} \ur{ضروری} \ur{ہے۔}
\end{quote}

\begin{quote}
\emph{musāfiroṅ kī hifāzat zarūrī hai.}
\end{quote}

\begin{quote}
Keeping passengers safe matters.
\end{quote}

Feminine.

\begin{cousinweb}
\textbf{\ur{حفاظت}} (\emph{hifāzat}) is usually traced, through Persian, to Arabic \emph{ḥifāẓa}, \textquotedblleft{}guarding, keeping safe\textquotedblright{}; \textbf{\ur{محفوظ}}, \textquotedblleft{}safe\textquotedblright{}, belongs to the same family.
\end{cousinweb}

\begin{grammarlens}[title={What you've built}]
One more word for getting around.
\end{grammarlens}

\subsection*{Your turn}
\begin{itemize}
\raggedright
  \item \emph{Say it:} \emph{hifāzat}
  \item \emph{Say it:} \emph{hifāzat}, once more
  \item \emph{Say it:} \emph{musāfiroṅ kī hifāzat zarūrī hai.} — Keeping passengers safe matters.
  \item \emph{Recall:} say the Urdu for to tighten, then the Urdu for postponed, then say \emph{hifāzat} again
\end{itemize}

\begin{tcolorbox}[breakable,colback=teal!4,colframe=teal!35!black,title={Before you move on}]
What does \ur{حفاظت} mean? (\textquotedblleft{}Protection; safety\textquotedblright{}.) Say it once more.
\end{tcolorbox}
\section[\texorpdfstring{\ur{تلاشی} (talāshī)}{تلاشی (talāshī)}]{\ur{تلاشی} (talāshī) — a security search; a search of bags}

\label{lesson:UR-C399-talashi}



\begin{quote}
Before the new one: say the Urdu for postponed, then the Urdu for protection.
\end{quote}

\subsection*{You'll want to know: \ur{تلاشی}}
\textbf{\ur{تلاشی}} — \emph{talāshī} — \textquotedblleft{}a security search; a search of bags\textquotedblright{}.

\begin{quote}
\ur{حفاظت} \ur{کے} \ur{لیے} \ur{ہمارے} \ur{تھیلوں} \ur{کی} \ur{تلاشی} \ur{ہوئی۔}
\end{quote}

\begin{quote}
\emph{hifāzat ke liye hamāre thailoṅ kī talāshī huī.}
\end{quote}

\begin{quote}
Our bags were searched for safety.
\end{quote}

Feminine.

\begin{cousinweb}
is built on \emph{talāsh}, \textquotedblleft{}a search\textquotedblright{}, usually traced to Persian \emph{talāsh} and perhaps further to Turkic, with the ending \emph{-ī}: \textbf{\ur{تلاش} \ur{کرنا}} is to look for something, \textbf{\ur{تلاشی}} the search of a bag or a person.
\end{cousinweb}

\begin{grammarlens}[title={What you've built}]
One more word for getting around.
\end{grammarlens}

\subsection*{Your turn}
\begin{itemize}
\raggedright
  \item \emph{Say it:} \emph{talāshī}
  \item \emph{Say it:} \emph{talāshī}, once more
  \item \emph{Say it:} \emph{hifāzat ke liye hamāre thailoṅ kī talāshī huī.} — Our bags were searched for safety.
  \item \emph{Recall:} say the Urdu for postponed, then the Urdu for protection, then say \emph{talāshī} again
\end{itemize}

\begin{tcolorbox}[breakable,colback=teal!4,colframe=teal!35!black,title={Before you move on}]
What does \ur{تلاشی} mean? (\textquotedblleft{}A security search; a search of bags\textquotedblright{}.) Say it once more.
\end{tcolorbox}
\section[\texorpdfstring{\ur{کسٹم} (kasṭam)}{کسٹم (kasṭam)}]{\ur{کسٹم} (kasṭam) — customs (at a border)}

\label{lesson:UR-C399-kastam}



\begin{quote}
Before the new one: say the Urdu for protection, then the Urdu for a security search.
\end{quote}

\subsection*{You'll want to know: \ur{کسٹم}}
\textbf{\ur{کسٹم}} — \emph{kasṭam} — \textquotedblleft{}customs (at a border)\textquotedblright{}.

\begin{quote}
\ur{کسٹم} \ur{والے} \ur{نے} \ur{میرا} \ur{صندوق} \ur{کھولا۔}
\end{quote}

\begin{quote}
\emph{kasṭam vāle ne merā sandūq kholā.}
\end{quote}

\begin{quote}
The customs officer opened my trunk.
\end{quote}

Masculine.

\begin{cousinweb}
\textbf{\ur{کسٹم}} (\emph{kasṭam}) is borrowed from English \emph{customs}.
\end{cousinweb}

\begin{grammarlens}[title={What you've built}]
One more word for getting around.
\end{grammarlens}

\subsection*{Your turn}
\begin{itemize}
\raggedright
  \item \emph{Say it:} \emph{kasṭam}
  \item \emph{Say it:} \emph{kasṭam}, once more
  \item \emph{Say it:} \emph{kasṭam vāle ne merā sandūq kholā.} — The customs officer opened my trunk.
  \item \emph{Recall:} say the Urdu for protection, then the Urdu for a security search, then say \emph{kasṭam} again
\end{itemize}

\begin{tcolorbox}[breakable,colback=teal!4,colframe=teal!35!black,title={Before you move on}]
What does \ur{کسٹم} mean? (\textquotedblleft{}Customs (at a border)\textquotedblright{}.) Say it once more.
\end{tcolorbox}
\section[\texorpdfstring{\ur{ممنوع} (mamnū)}{ممنوع (mamnū)}]{\ur{ممنوع} (mamnū) — forbidden; not allowed}

\label{lesson:UR-C399-mamnu}



\begin{quote}
Before the new one: say the Urdu for a security search, then the Urdu for customs (at a border).
\end{quote}

\subsection*{You'll want to know: \ur{ممنوع}}
\textbf{\ur{ممنوع}} — \emph{mamnū} — \textquotedblleft{}forbidden; not allowed\textquotedblright{}.

\begin{quote}
\ur{سرحد} \ur{پر} \ur{یہ} \ur{چیز} \ur{ممنوع} \ur{ہے۔}
\end{quote}

\begin{quote}
\emph{sarhad par yih chīz mamnū hai.}
\end{quote}

\begin{quote}
This item is not allowed across the border.
\end{quote}

\begin{cousinweb}
\textbf{\ur{ممنوع}} (\emph{mamnū}) is usually traced, through Persian, to Arabic \emph{mamnūʿ}, \textquotedblleft{}forbidden\textquotedblright{}; signs at airports and borders use it.
\end{cousinweb}

\begin{grammarlens}[title={What you've built}]
One more word for getting around.
\end{grammarlens}

\subsection*{Your turn}
\begin{itemize}
\raggedright
  \item \emph{Say it:} \emph{mamnū}
  \item \emph{Say it:} \emph{mamnū}, once more
  \item \emph{Say it:} \emph{sarhad par yih chīz mamnū hai.} — This item is not allowed across the border.
  \item \emph{Recall:} say the Urdu for a security search, then the Urdu for customs (at a border), then say \emph{mamnū} again
\end{itemize}

\begin{tcolorbox}[breakable,colback=teal!4,colframe=teal!35!black,title={Before you move on}]
What does \ur{ممنوع} mean? (\textquotedblleft{}Forbidden; not allowed\textquotedblright{}.) Say it once more.
\end{tcolorbox}
\section[\texorpdfstring{\ur{چڑھانا} (chaṛhānā)}{چڑھانا (chaṛhānā)}]{\ur{چڑھانا} (chaṛhānā) — to put on board; to put up}

\label{lesson:UR-C399-charhana}



\begin{quote}
Before the new one: say the Urdu for customs (at a border), then the Urdu for forbidden.
\end{quote}

\subsection*{You'll want to know: \ur{چڑھانا}}
\textbf{\ur{چڑھانا}} — \emph{chaṛhānā} — \textquotedblleft{}to put on board; to put up\textquotedblright{}.

\begin{quote}
\ur{مزدوروں} \ur{نے} \ur{صندوق} \ur{ٹرک} \ur{پر} \ur{چڑھائے۔}
\end{quote}

\begin{quote}
\emph{mazdūroṅ ne sandūq ṭrak par chaṛhāe.}
\end{quote}

\begin{quote}
The labourers put the trunks up onto the lorry.
\end{quote}

Past, with \ur{نے} (he or she …): \ur{اس} \ur{نے} \ur{سامان} \ur{اوپر} \ur{چڑھایا}, \emph{us ne sāmān ūpar chaṛhāyā}.

\begin{cousinweb}
\textbf{\ur{چڑھانا}} (\emph{chaṛhānā}) is built on \textbf{\ur{چڑھنا}}, \textquotedblleft{}to climb\textquotedblright{}, which comes from \emph{caḍhati}, \textquotedblleft{}rises\textquotedblright{}, a form CDIAL reconstructs (CDIAL 4578): this is the form for raising something else.
\end{cousinweb}

\begin{grammarlens}[title={What you've built}]
One more word for getting around.
\end{grammarlens}

\subsection*{Your turn}
\begin{itemize}
\raggedright
  \item \emph{Say it:} \emph{chaṛhānā}
  \item \emph{Say it:} \emph{chaṛhānā}, once more
  \item \emph{Say it:} \emph{mazdūroṅ ne sandūq ṭrak par chaṛhāe.} — The labourers put the trunks up onto the lorry.
  \item \emph{Recall:} say the Urdu for customs (at a border), then the Urdu for forbidden, then say \emph{chaṛhānā} again
\end{itemize}

\begin{tcolorbox}[breakable,colback=teal!4,colframe=teal!35!black,title={Before you move on}]
What does \ur{چڑھانا} mean? (\textquotedblleft{}To put on board; to put up\textquotedblright{}.) Say it once more.
\end{tcolorbox}
